% !TEX root = ./main.tex

% 模板默认 zh+en 双扉页；双面 + openright 会在英文扉页前插入空白页。
% 仅保留中文扉页时改为只含 zh（须在 \sjtusetup 之前执行）。
\ExplSyntaxOn
\clist_gset:Nn \g__sjtu_lang_clist { zh }
\ExplSyntaxOff

\sjtusetup{
  %
  %******************************
  % 注意：
  %   1. 配置里面不要出现空行
  %   2. 不需要的配置信息可以删除
  %******************************
  %
  % 信息录入
  %
  info = {%
    %
    % 标题
    %
    zh / title           = {基于 Flink 的动态 Kafka-Topic 发现与订阅机制研究},
    en / title           = {A Sample Document for \LaTeX-based SJTU Thesis Template},
    %
    % 标题页标题
    %   可使用“\\”命令手动控制换行
    %
    % zh / display-title   = {上海交通大学学位论文\\ \LaTeX{} 模板示例文档},
    % en / display-title   = {A Sample Document \\ for \LaTeX-based SJTU Thesis Template},
    %
    % 关键词
    %
    zh / keywords        = {Flink, Kafka, 动态 Topic, Sink, 流处理, 路由映射},
    en / keywords        = {SJTU, master thesis, XeTeX/LaTeX template},
    %
    % 姓名
    %
    zh / author          = {胡文杰},
    en / author          = {Mo Mo},
    %
    % 学号
    %
    id              = {522031910511},
    %
    % 指导教师
    %
    zh / supervisor      = {任锐},
    en / supervisor      = {Prof.\ Mou Mou},
    %
    % 副指导教师
    %
    % zh / assoc-supervisor  = {某某教授},
    % en / assoc-supervisor  = {Prof.\ Uom Uom},
    %
    % 联合导师
    % zh / co-supervisor     = {某某工程师},
    % en / co-supervisor     = {Engr.\ Oum Oum},
    %
    % 所属院系
    %
    zh / department      = {计算机学院},
    en / department      = {Depart of XXX},
    %
    % 专业
    %
    zh / major           = {软件工程},
    en / major           = {A Very Important Major},
    %
    % 学位
    %   除交叉学科门类外，各一级学科按所属学科门类申请学位
    %   专业学位类别按专业名称申请学位
    %   本科生不需要填写
    %
    zh / degree          = {工学学士},
    en / degree          = {Bachelor of Engineering},
    %
    % 答辩日期
    %   使用 ISO 格式 (yyyy-mm-dd)；默认为当前时间
    %
    % date                 = {2023-05-18},
    %
    % 标题页显示日期
    %   覆盖对应标题页的日期显示，原样输出
    %
    % zh / display-date    = {2023 年 5 月},
    %
    % 资助基金
    %
    % zh / fund  = {
    %                {国家 973 项目 (No.\ 2025CB000000)},
    %                {国家自然科学基金 (No.\ 81120250000)},
    %              },
    % en / fund  = {
    %                {National Basic Research Program of China (Grant No.\ 2025CB000000)},
    %                {National Natural Science Foundation of China (Grant No.\ 81120250000)},
    %              },
  },
  %
  % 风格设置
  %
  style = {%
    %
    % 关键词首行悬挂
    %
    % keywords-format = hang,
  },
  %
  % 名称设置
  %
  name = {
    % bib             = {References},
    % ack             = {谢\hspace{\ccwd}辞},
    % achv            = {攻读学位期间完成的论文},
  },
}

% 使用 BibLaTeX 处理参考文献
%   biblatex-gb7714-2015 常用选项
%     gbnamefmt=lowercase     姓名大小写由输入信息确定
%     gbpub=false             禁用出版信息缺失处理
\usepackage[backend=biber,style=gb7714-2015]{biblatex}
% 文献表字体
\renewcommand{\bibfont}{\zihao{5}\setbaselineskip{16bp}}
% 文献表条目间的间距
\setlength{\bibitemsep}{3bp plus 1pt}
% 导入参考文献数据库
\addbibresource{refs.bib}

% 脚注格式
\usepackage[perpage,bottom,hang]{footmisc}

% Provide placeholder images like `example-image-a` / `example-image-b`
\usepackage{mwe}

% 定义图片文件目录与扩展名
\graphicspath{{figures/}}
\DeclareGraphicsExtensions{.pdf,.eps,.png,.jpg,.jpeg}

% 颜色宏包
\usepackage{xcolor} 

% 确定浮动对象的位置，可以使用 [H]，强制将浮动对象放到这里（可能效果很差）
% \usepackage{float}

% 固定宽度的表格
% \usepackage{tabularx}

% 使用三线表：toprule，midrule，bottomrule。
\usepackage{booktabs}

% 表格中支持跨行
\usepackage{multirow}

% 表格中数字按小数点对齐
\usepackage{dcolumn}
\newcolumntype{d}[1]{D{.}{.}{#1}}

% 使用长表格
\usepackage{longtable}

% 附带脚注的表格
\usepackage{threeparttable}

% 附带脚注的长表格
\usepackage{threeparttablex}

% 算法环境宏包
\usepackage[ruled,vlined,linesnumbered]{algorithm2e}
% \usepackage{algorithm, algorithmicx, algpseudocode}

% 代码环境宏包
\usepackage{listings}
\lstdefinestyle{lstStyleCode}{%
  aboveskip         = \medskipamount,
  belowskip         = \medskipamount,
  basicstyle        = \ttfamily\zihao{-5}\setbaselineskip{12bp},
  commentstyle      = \slshape\color{black!60},
  stringstyle       = \color{green!40!black!100},
  keywordstyle      = \bfseries\color{blue!50!black},
  extendedchars     = false,
  upquote           = true,
  tabsize           = 2,
  showstringspaces  = false,
  xleftmargin       = 1em,
  xrightmargin      = 1em,
  breaklines        = true,
  framexleftmargin  = 1em,
  framexrightmargin = 1em,
  backgroundcolor   = \color{gray!10},
  columns           = flexible,
  keepspaces        = true,
  % texcl             = true,  % 代码中的注释部分将会被视为 LaTeX 正文模式内容
  % mathescape        = true,  % 代码中$..$的部分将会被视为 LaTeX 数学模式内容
}
\lstnewenvironment{codeblock}[1][]{%
  \lstset{style=lstStyleCode,#1}}{}

% 直立体数学符号
\providecommand{\dd}{\mathop{}\!\mathrm{d}}
\providecommand{\ee}{\mathrm{e}}
\providecommand{\ii}{\mathrm{i}}
\providecommand{\jj}{\mathrm{j}}

% 国际单位制宏包
\usepackage{siunitx}

% 定理环境宏包
\usepackage{amsthm}
% \usepackage{ntheorem}

% 绘图宏包
\usepackage{tikz}
\usetikzlibrary{arrows.meta, shapes.geometric}

% 数据图表宏包
\usepackage{pgfplots}
\pgfplotsset{compat=newest}

% 一些文档中用到的 logo
\usepackage{hologo}
\providecommand{\XeTeX}{\hologo{XeTeX}}
\providecommand{\BibLaTeX}{\textsc{Bib}\LaTeX}

% 借用 ltxdoc 里面的几个命令方便写文档
\DeclareRobustCommand\cs[1]{\texttt{\char`\\#1}}
\providecommand\pkg[1]{{\sffamily#1}}

% hyperref 宏包在最后调用
% hidelinks：不在页面上绘制链接彩色边框（目录等处否则常显示为红框）
\usepackage[hidelinks]{hyperref}

% E-mail
\providecommand{\email}[1]{\href{mailto:#1}{\urlstyle{tt}\nolinkurl{#1}}}

% Give TeX a little extra room to break long mixed Chinese/identifier lines.
\setlength{\emergencystretch}{3em}

% 双面 openright 时，模板在 \cleardoublepage 插入的空白页上使用 \thispagestyle{empty}，
% 故无页眉页码；改为与普通页一致（仍无正文，仅恢复页眉页码）
\makeatletter
\RenewDocumentCommand \cleardoublepage { }
  {%
    \clearpage
    \if@twoside
      \ifodd\c@page
      \else
        \hbox{}%
        \newpage
        \if@twocolumn
          \hbox{}%
          \newpage
        \fi
      \fi
    \fi
  }
\makeatother
